% 天王星（综述）
% license CCBYSA3
% type Wiki

本文根据 CC-BY-SA 协议转载翻译自维基百科\href{https://en.wikipedia.org/wiki/Uranus}{相关文章}。
\begin{figure}[ht]
\centering
\includegraphics[width=8cm]{./figures/fb24dd1c06777ad6.png}
\caption{由旅行者 2 号拍摄的天王星真彩图像\(^\text{[a]}\)。其苍白而柔和的外观源于覆盖在云层之上的一层薄雾。} \label{fig_Uranus_1}
\end{figure}
天王星是距离太阳的第七颗行星。它是一颗呈青色的气态冰巨行星。该行星的大部分由处于超临界相状态的水、氨和甲烷构成，天文学上将这些物质称为“冰”或挥发物。该行星的大气具有复杂的分层云结构，并在所有太阳系行星中具有最低的最低温度（49 K（−224 °C; −371 °F））。它具有 82.23° 的显著自转轴倾角，自转周期为 17 小时 14 分钟，且为逆行自转。这意味着在其围绕太阳完成一个为期 84 个地球年的公转周期期间，其两极会经历约 42 年的连续日照，随后是约 42 年的连续黑暗。

在太阳系行星中，天王星直径排名第三、质量排名第四。根据当前模型，在其挥发性地幔层内部存在一个岩质核心，外围包覆着厚厚的氢和氦大气。在其高层大气中检测到痕量的碳氢化合物（被认为通过水解作用产生）以及一氧化碳和二氧化碳（被认为源自彗星）。天王星大气中存在许多尚未解释的气候现象，例如其峰值风速 900 km/h（560 mph）^{\text{[24]}}、极冠的变化以及其不规则的云层结构。该行星的内部热量相较其他巨行星异常低，其原因尚不清楚。

与其他巨行星一样，天王星拥有环系统、磁层以及众多天然卫星。其极为暗淡的环系统仅反射约 2\% 的入射光。天王星的 29 颗天然卫星中包含 19 颗已知的规则卫星，其中 14 颗为小型内侧卫星。在更外侧的是该行星五颗较大的主要卫星：米兰达（Miranda）、艾丽尔（Ariel）、安布里尔（Umbriel）、泰坦妮亚（Titania）和奥伯龙（Oberon）。在距离天王星更远的轨道上则有十颗已知的不规则卫星。该行星的磁层高度不对称，且含有大量带电粒子，这可能是其环和卫星变暗的原因。

天王星肉眼可见，但亮度极低，并且相对于背景恒星移动非常缓慢，因此直到 1781 年被威廉·赫歇尔首次观测到之前，它从未被归类为一颗行星。在其发现约七十年后，人们达成共识，将其命名为希腊原始神祇之一乌拉诺斯（Uranus，Ouranos）。截至 2025 年，它仅被探测器访问过一次，即 1986 年旅行者二号（Voyager 2）飞掠该行星^{\text{[25]}}。尽管如今可通过望远镜分辨并观测其细节，但科研界强烈希望再次造访该行星——这一点体现于《行星科学十年调查》将拟议的“天王星轨道器与探测器任务”列为 2023–2032 年的最高优先级任务之一，以及中国国家航天局提出利用天问四号的子探测器飞掠天王星的计划^{\text{[26]}}。
\subsection{历史}
1781 年 3 月 13 日（天王星被发现之日），天王星的位置（以叉号标示）与古典行星一样，天王星肉眼可见，但由于其亮度极低且公转缓慢，古代观测者从未将其识别为一颗行星^{\text{[27]}}。威廉·赫歇尔于 1781 年 3 月 13 日首次观测到天王星，由此发现了这颗行星，这也是人类历史上首次扩展太阳系已知边界，并使天王星成为第一颗借助望远镜而被确认的行星。天王星的发现实际上使当时已知太阳系的尺度扩大了一倍，因为天王星距离太阳的平均距离约为土星的两倍。
\subsubsection{发现}
在其被确认是一颗行星之前，天王星曾被多次观测到，但通常被误认为是一颗恒星。显然，喜帕恰斯在公元前 128 年观测过它，他测量了恒星的位置以编制星表，该星表后来被纳入托勒密的《天文大成》。该星表给出了处于室女座的四颗恒星的位置，它们构成一个四边形。其中有一颗恒星实际上并不存在，而在公元前 128 年 4 月，天王星正位于该位置^{\text{[28]}}。最早的可靠观测记录是在 1690 年，当时约翰·弗拉姆斯蒂德至少六次观测到它，并将其编入星表为金牛座 34 号星。詹姆斯·布拉德雷在 1748 年、1750 年和 1753 年三次观测到它。托比亚斯·迈耶于 1756 年观测到它一次。法国天文学家皮埃尔·夏尔·勒莫尼耶在 1750 年至 1769 年期间至少十二次观测到天王星，其中包括连续四个夜晚的观测^{\text{[29]}}。

威廉·赫歇尔于 1781 年 3 月 13 日在英国萨默塞特郡巴斯市新金街 19 号自家花园中（现为赫歇尔天文学博物馆）^{\text{[30]}}观测到天王星，并在 1781 年 4 月 26 日最初将其报告为一颗彗星^{\text{[31]}}。赫歇尔使用自制的 6.2 英寸反射式望远镜，“从事一系列关于恒星视差的观测”^{\text{[32][33]}}。

赫歇尔在日志中写道：“在金牛座 ζ 附近的四分仪区域中……要么是一颗‘星云状恒星’，要么或许是一颗彗星。”^{\text{[34]}} 3 月 17 日，他又记道：“我寻找那颗彗星或星云状恒星，发现它是一颗彗星，因为它改变了位置。”^{\text{[35]}} 当他向皇家学会提交他的发现时，他继续坚称自己发现了一颗彗星，但也在言语间将其与一颗行星作了比较^{\text{[32]}}：

“我第一次看到这颗彗星时所用的放大倍率是 227 倍。根据经验我知道，固定恒星的视直径在更高放大倍率下并不会按比例放大，而行星会；因此我将放大倍率调至 460 倍和 932 倍，发现彗星的直径确实随着放大倍率按比例增大，正如在假设它不是固定恒星时所应当出现的情况，而我用来比较的那些恒星的直径并未按相同比例增大。此外，这颗彗星在远超出其亮度所允许的放大倍率下，显得朦胧且轮廓不清，而恒星却保持了那种我从数以千计的观测中所熟知的光亮和清晰度。随后的一切证明了我的推测是正确的，这正是我们最近观测到的那颗彗星。”

赫歇尔将他的发现通知了皇家天文学家内维尔·马斯克林，而马斯克林在 1781 年 4 月 23 日给他回信道出困惑：“我不知道该称它为何物。它有可能是一颗沿着近乎圆形轨道绕太阳运行的常规行星，也有可能是一颗沿着高度偏心椭圆运行的彗星。我尚未看到它有任何彗发或彗尾。”^{\text{[36]}}
